\subsection{The mass of a maximal integral Hermitian lattice}
\label{sec:maximal-hermitian-mass}

The genus count in the degenerate return sum requires the mass of a
maximal integral lattice, which need not be unimodular.  The following
formula supplies that input.  It follows from the classical Hermitian
mass formula; the local argument below identifies every correction
factor, including at primes above \(2\).

Let \(E/F\) be a CM quadratic extension, unramified at every finite
place, with \(F\) totally real of degree \(f\).  Write \(\mathcal O_E\)
and \(\mathcal O_F\) for their rings of integers, \(D_F\) for the
absolute discriminant of \(F\), and \(\eta\) for the quadratic Hecke
character associated with \(E/F\).  Thus, at a finite prime
\(\mathfrak p\) of \(F\), \(\eta(\mathfrak p)=1\) if the prime splits
and \(\eta(\mathfrak p)=-1\) if it is inert.  All \(L\)-functions below
are finite Hecke \(L\)-functions, without their archimedean factors.

Let \(V\) be an \(n\)-dimensional Hermitian \(E\)-space, \(n\geq1\),
positive definite at every real place of \(F\).  An integral lattice
is a full, finitely generated \(\mathcal O_E\)-submodule \(L\subset V\)
such that \(h(L,L)\subset\mathcal O_E\).  Its dual is
\[
 L^\#=\{x\in V:h(x,L)\subset\mathcal O_E\}.
\]
The lattice is maximal integral if it has no proper integral
overlattice.  Its genus consists of the lattices isometric to \(L\)
at every finite place; write
\[
 \operatorname{Mass}(L)
   =\sum_{[M]\in\operatorname{gen}(L)}
                         \frac1{|\operatorname{Aut}_{\mathcal O_E,h}(M)|}.
\]
All automorphisms here are \(E\)-linear isometries.  In particular, no
determinant-one condition and no extra quotient by scalar
automorphisms is imposed.

\begin{theorem}[Maximal integral Hermitian mass]
\label{thm:maximal-hermitian-mass}
For fixed \(E\) and \(n\), set
\[
 C(E,n)=
 2D_F^{n^2/2}
 \left(\prod_{a=1}^{n}\frac{(a-1)!}{(2\pi)^a}\right)^f
 \prod_{a=1}^{n}L_F(a,\eta^a).
 \tag{MM1}
\]
This is a positive real number independent of \(V\) and \(L\).
If \(L\) is maximal integral, all primes of \(\mathcal O_E\) at which
\(L\) is not unimodular are inert over \(F\).  There are finitely many
such primes.  If \(B(L)\) denotes this set and
\(q_P=\sqrt{N_{E/\mathbf Q}P}\), then
\[
 \operatorname{Mass}(L)
   =C(E,n)\prod_{P\in B(L)}
               \frac{q_P^{\,n}-(-1)^n}{q_P+1}.
 \tag{MM2}
\]
For \(n=1\), every factor in this finite product is \(1\).
\end{theorem}

\begin{proof}
We first determine the local lattices.  Maximality is a local
condition: if a completion has a proper integral overlattice, its
quotient by the original completion is a nonzero submodule of the
corresponding primary part of \(L^\#/L\).  Taking the inverse image
of that submodule in \(L^\#\) gives a global proper overlattice.
Its integrality can be tested at every finite place.  Conversely,
a proper global integral overlattice remains proper at a prime
dividing its finite index.

At a split place \(v\), write
\(E_v=F_v\times F_v\); conjugation exchanges the factors.  The local
lattice decomposes as \(L_+\oplus L_-\), and the Hermitian form is
determined by a perfect pairing between the two corresponding
vector spaces.  Integrality says \(L_-\subset L_+^\vee\).
Replacing \(L_-\) by \(L_+^\vee\) preserves integrality, so
maximality gives equality.  The pairing on the two modules is then
perfect, and \(L_v\) is unimodular.

Now let \(v\) be inert.  Put \(A=\mathcal O_{F_v}\),
\(B=\mathcal O_{E_v}\), and choose a uniformizer \(\pi\in A\).
The extension \(B/A\) is unramified of degree two.  An integral
Hermitian \(B\)-lattice has an orthogonal basis with Gram matrix
\[
 \operatorname{diag}(\pi^{a_1}u_1,\ldots,\pi^{a_n}u_n),
 \qquad a_i\geq0,\quad u_i\in A^\times.
 \tag{MM3}
\]
For completeness, let \(s\) be the least valuation of a Gram-matrix
entry.  If a diagonal entry has valuation \(s\), its vector can
be split off.  Otherwise, take two basis vectors whose off-diagonal
pairing has valuation \(s\).  The trace pairing of the quadratic
residue-field extension is nondegenerate, so a linear combination
of those vectors has norm of valuation \(s\).  Division by that
norm gives an integral orthogonal projection onto its line.
Induction gives (MM3).  This argument uses the residue-field trace,
which is surjective also in characteristic \(2\); it does not
divide by \(2\).

If any \(a_i\geq2\), adjoining \(\pi^{-1}e_i\) gives a proper
integral overlattice.  Hence \(a_i\in\{0,1\}\).  If two exponents
are \(1\), the corresponding binary form after division by \(\pi\)
has a nonzero isotropic vector over the residue field.  Indeed,
surjectivity of
\(\mathbf F_{q^2}^{\times}\to\mathbf F_q^\times\),
\(z\mapsto z\overline z\), solves
\(u_1+u_2z\overline z=0\).  Lifting such a vector to a primitive
vector \(w\in Be_i+Be_j\) gives
\[
 h(w,L_v)\subset\pi B,\qquad h(w,w)\in\pi^2A.
\]
Adjoining \(\pi^{-1}w\) again gives a proper integral overlattice.
Thus a maximal lattice is either unimodular or has the form
\[
 L_v=L_{v,0}\perp Be,\qquad
 L_{v,0}\text{ unimodular of rank }n-1,\qquad h(e,e)=\pi u
 \quad(u\in A^\times).
 \tag{MM4}
\]
Conversely, (MM4) is maximal.  Its dual quotient is one-dimensional
over \(B/\pi B\).  Every nonzero element of that quotient has a
representative whose norm has valuation \(-1\); consequently it
cannot lie in an integral overlattice.  This also proves
\[
 |L_v^\#/L_v|=q^2
 \quad\text{at a nonunimodular place,}
 \tag{MM5}
\]
where \(q=N_{F/\mathbf Q}v\).  At this inert place \(q_P=q\).
The quotient \(L^\#/L\) is finite, so only finitely many such places
occur.

We use the classical Siegel mass formula for Hermitian lattices
in the normalization
\[
 \operatorname{Mass}(L)=
 2D_F^{n^2/2}
 \left(\prod_{a=1}^{n}\frac{(a-1)!}{(2\pi)^a}\right)^f
 |L^\#/L|^{n/2}\prod_v\beta(L_v)^{-1}.
 \tag{MM6}
\]
Here the relative discriminant of \(E/F\) is \(1\), and
\[
 \beta(L_v)=
 \lim_{N\longrightarrow\infty}
 q^{-Nn^2}
 \#\{g\in\operatorname{End}_{B/\pi^N B}(L_v/\pi^NL_v):
                    h(gx,gy)=h(x,y)\bmod\pi^N\}.
 \tag{MM7}
\]
The split version uses \(B=A\times A\).  Formula (MM6), the
stabilized-density convention (MM7), and the Jordan-block density
formula are the classical mass theorem as stated in
\cite[Theorems 4.2.3--4.2.4 and Lemma 4.2.5]{kirschmer-herm};
its proof is attributed
there to Gan--Yu, Section 10, and the local formula to their
Theorem 7.3.  In the present setting every finite prime is a good
prime in that theorem: the excluded unitary case is a ramified
quadratic extension in residue characteristic \(2\), which is
absent here.  The base field \(F_v/\mathbf Q_2\) need not be
unramified.

Define
\[
 b_t(q,\epsilon)=\prod_{a=1}^{t}(1-\epsilon^a q^{-a}),
 \qquad b_0(q,\epsilon)=1.
 \tag{MM8}
\]
At a unimodular place the local factor is
\(\beta(L_v)=b_n(q,\eta(v))\).  At an inert place of type (MM4),
the two Jordan blocks have ranks \(n-1,1\), exponents \(0,1\),
and reductive quotients \(U_{n-1}\) and \(U_1\).  In the cited
local density formula their power of \(q\) is \(q^1\): the first
block contributes exponent zero, the second contributes one,
and the cross term has the first block's exponent zero.  Using
\[
 |U_t(\mathbf F_q)|
   =q^{t(t-1)/2}\prod_{a=1}^{t}(q^a-(-1)^a)
\]
therefore gives
\[
 \beta(L_v)
 =q\,b_{n-1}(q,-1)b_1(q,-1)
 =(q+1)b_{n-1}(q,-1).
 \tag{MM9}
\]
For \(n=1\), this is the single exponent-one rank-one block and
the same formula holds with \(b_0=1\).  Notice that the density
is a stabilized congruence count.  At a nonunimodular place one
cannot replace it by the count modulo \(\pi\) without checking
stabilization.

Relative to a unimodular local factor, (MM5)--(MM9) change the
right side of (MM6) by exactly
\[
 \frac{q^n b_n(q,-1)}{(q+1)b_{n-1}(q,-1)}
                =\frac{q^n-(-1)^n}{q+1}.
 \tag{MM10}
\]
At all other places the change is \(1\).  The remaining standard
Euler factors are
\[
 \prod_v b_n(q_v,\eta(v))^{-1}
                    =\prod_{a=1}^{n}L_F(a,\eta^a),
\]
with the \(a=1\) factor understood by its Hecke \(L\)-value, as
in the mass formula.  The factor \(L_F(1,\eta)\) is positive
by the quadratic relative class number formula; all other
factors have positive convergent Euler products.  This gives
(MM1)--(MM2), uniformly for every positive definite Hermitian
space of rank \(n\).
\end{proof}

\paragraph{The exact input supplied to the return count.}
In the original source, the required maximal-lattice mass statement
quantifies first over \(E\) and \(n\), then over all positive definite
rank-\(n\) Hermitian spaces and their maximal integral lattices.
Its product runs over prime ideals \(P\) of \(\mathcal O_E\) at
which the lattice is not unimodular, with
\(q_P=\lfloor\sqrt{N P}\rfloor\).
The local proof above shows that each such prime is inert and
\(N P=q_P^2\); it also identifies this product exactly with (MM2).
Thus Theorem~\ref{thm:maximal-hermitian-mass} supplies that entire
mass statement, with \(C(E,n)\) chosen before either the space or
the lattice varies.\cite{maximal-hermitian-mass-source}

This closes the mass input, but not the fixed-label return count.
That count still requires uniform genus averages of primitive
sublattices, the counts of isometries with prescribed fixed space,
and the original row-to-lattice correspondence.  None of those
estimates follows from the mass formula alone.

\begin{corollary}[The mass of the original cyclic-discriminant genus]
\label{cor:cyclic-discriminant-mass}
Keep \(E/F\) as above.  Let \(T\) be a rational prime that splits
completely in \(E\), and let \(L\) be an integral positive definite
Hermitian lattice of rank \(n\geq1\) such that
\[
 L^\#/L\simeq\mathcal O_E/T\mathcal O_E
                     \quad\text{as }\mathcal O_E\text{-modules}.
 \tag{MM11}
\]
Then
\[
 \operatorname{Mass}(L)
  =C(E,n)\left(\frac{T^n-1}{T-1}\right)^f.
 \tag{MM12}
\]
If \(I\) is any unimodular positive definite Hermitian lattice of
rank \(n\), then \(C(E,n)=\operatorname{Mass}(I)\).  Moreover, any
two lattices satisfying (MM11) in the same Hermitian space belong
to the same genus.
\end{corollary}

\begin{proof}
Away from \(T\), (MM11) makes \(L_v\) unimodular.  At a place over
\(T\), both components of the split algebra have a one-dimensional
dual quotient over the residue field.  The elementary divisor
theorem applied to the pairing between \(L_+\) and \(L_-\) gives
one exponent \(1\) and \(n-1\) exponents \(0\).  The two Jordan
blocks now have reductive quotients \(GL_{n-1}\) and \(GL_1\).
The same local density formula as above gives
\[
 \beta(L_v)=T\,b_{n-1}(T,1)b_1(T,1)
                         =(T-1)b_{n-1}(T,1).
\]
The local discriminant quotient has size \(T^2\), so the
correction to the standard local factor is
\[
 \frac{T^n b_n(T,1)}
             {(T-1)b_{n-1}(T,1)}
                         =\frac{T^n-1}{T-1}.
\]
There are exactly \(f\) primes of \(F\) over \(T\).  Substitution
in (MM6) proves (MM12), and the same formula with all local
exponents zero proves the assertion about \(I\).

The elementary divisors determine the split local lattice up to
isometry.  At an inert place, a unimodular lattice is also unique
up to isometry: the integral diagonalization used in (MM3) has
only unit coefficients, and every unit of \(A\) is a norm from
the unramified quadratic extension \(B\).  To see the last fact
without excluding residue characteristic \(2\), first use
surjectivity of the residue-field norm, then lift one power of
\(\pi\) at a time using
\[
 N(1+\pi^a z)\equiv
          1+\pi^a\operatorname{Tr}(\bar z)\pmod{\pi^{a+1}}
                    \qquad(a\geq1).
\]
The residue-field trace is onto, and completeness gives the
desired unit norm.  Thus two lattices satisfying (MM11) are
isometric at every finite place, proving the genus assertion.
\end{proof}

The factor in (MM12) is the original normalization in the
primitive-sublattice genus mean.  In particular its size and its
dependence on the degree \(f\) are not replaced by an unspecified
constant.\cite{maximal-hermitian-mass-source}

\begin{theorem}[Mass at a fixed discriminant in a fixed space]
\label{thm:integral-lattice-fixed-discriminant-mass}
Let \(V\) be a fixed positive definite Hermitian \(E\)-space of
rank \(a\geq1\).  Put \(d=[E:\mathbf Q]=2f\) and
\(\ell=\log_2d\).  For every integer \(n\geq1\),
\[
 \sum_{\substack{[L]\text{ integral in }V\\ |L^\#/L|=n}}
                   \frac1{|\operatorname{Aut}(L)|}
          \leq C(E,a)\,n^{(a+\ell)/2}.
 \tag{MM13}
\]
The sum is over isometry classes inside this fixed rational space;
it is not a sum over different rational Hermitian spaces.
\end{theorem}

\begin{proof}
All maximal integral lattices in \(V\) form one genus.  At a split
place they are unimodular.  At an inert place the two possibilities
(MM3)--(MM4) are distinguished by the parity of the valuation of
their determinants, which is invariant under an \(E_v\)-linear
isometry.  Within either possibility the unit coefficients are
norms, so the local maximal lattice is unique up to isometry.
The rational space \(V\) fixes the parity, hence fixes the genus.

Choose a maximal integral lattice \(M\).  Such a lattice exists:
start with an integral full lattice and take a maximal integral
intermediate module between it and its dual.  There are only
finitely many intermediate modules.  By (MM5),
\[
 |M^\#/M|=D^2,\qquad
 D=\prod_{v\text{ nonunimodular}}q_v\in\mathbf Z_{\geq1}.
\]
The number \(D\) is the same throughout the maximal genus.
Theorem~\ref{thm:maximal-hermitian-mass} gives
\[
 \operatorname{Mass}(M)\leq C(E,a)D^a,
 \tag{MM14}
\]
because \((q^a-(-1)^a)/(q+1)\leq q^a\) for \(q\geq2\).

Every integral lattice \(L\) is contained in a maximal integral
lattice \(M'\), and such an overlattice lies in \(L^\#\).
Write \(z=[M':L]\), where indices are cardinalities of finite
additive quotients.  The dual index equality gives
\[
 |L^\#/L|=|M'^\#/M'|\,[M':L]^2=D^2z^2.
 \tag{MM15}
\]
Consequently the sum in (MM13) is empty unless \(n=D^2z^2\)
for an integer \(z\geq1\).  When this happens, \(z\) is fixed.
There is no sum over maximal discriminants or over divisors
arising from different choices of \(V\).

We need an elementary count of \(\mathcal O_E\)-sublattices of
index \(z\) in a fixed rank-\(a\) lattice.  Over a discrete
valuation ring with residue cardinality \(Q\), let \(A_a(e)\)
be the number of submodules of colength \(e\) in a free module
of rank \(a\).  Hermite normal form gives
\[
 \sum_{e\geq0}A_a(e)X^e
                 =\prod_{i=0}^{a-1}(1-Q^iX)^{-1}.
\]
In particular
\[
 A_a(e)=A_{a-1}(e)+Q^{a-1}A_a(e-1),\qquad
                         A_a(e)\leq Q^{ae}.
 \tag{MM16}
\]
For the inequality, use induction on \(a+e\), with \(e=0\)
or \(a=1\) immediate.  When \(a,e\geq1\), the two inductive
terms divided by \(Q^{ae}\) are at most \(Q^{-e}\) and
\(Q^{-1}\), whose sum is at most \(1\).

For a nonzero integral ideal \(\mathfrak a\), the quotient
lengths at its primes determine an index ideal.  Localizing
the lattice, applying (MM16), and using the Chinese remainder
theorem shows that the number of sublattices with index ideal
\(\mathfrak a\) is at most \((N\mathfrak a)^a\).
This does not require the lattice to be globally free.

Let \(I_E(z)\) be the number of ideals of norm \(z\).
At a rational prime \(p\), at most \(d\) prime ideals lie
above \(p\), and their residue degrees are positive.  The
coefficient of \(X^e\) in their ideal-counting product is
at most
\[
 [X^e](1-X)^{-d}
          =\binom{e+d-1}{d-1}\leq d^e.
 \tag{MM17}
\]
The last inequality counts weak compositions of \(e\) by
forgetting the order in words of length \(e\) on \(d\) letters.
Multiplying over rational primes and using \(p\geq2\) gives
\[
 I_E(z)\leq d^{\Omega(z)}\leq z^\ell.
\]
Thus a rank-\(a\) lattice has at most \(z^{a+\ell}\)
sublattices of additive index \(z\).

To retain the weights, count pairs \(L\subset M'\), with
\(M'\) maximal and \([M':L]=z\), modulo isometries of \(V\).
Grouping first by \(M'\) and applying the finite
orbit--stabilizer formula gives pair mass at most
\(\operatorname{Mass}(M)z^{a+\ell}\).
Grouping instead by \(L\) gives its mass multiplied by the
number of maximal integral overlattices of \(L\).  That
number is finite and at least one.  The required sum is
therefore at most the pair mass.  Equations (MM14)--(MM15)
now give
\[
 \operatorname{Mass}(M)z^{a+\ell}
 \leq C(E,a)D^az^{a+\ell}
 \leq C(E,a)(Dz)^{a+\ell}
 =C(E,a)n^{(a+\ell)/2}.
\]
All sums are finite: the maximal genus has finite class
number, and each of its representatives has finitely many
sublattices of index \(z\).  This proves (MM13).
\end{proof}

If the reference rank-\(a\) lattice is the original orthogonal
unimodular lattice \(I_a\), then \(C(E,a)=\operatorname{Mass}(I_a)\).
Thus (MM13) is the precise mass bound used in the original
primitive-sublattice genus count, including its
\(\log_2[E:\mathbf Q]\) exponent.
